% Clase 25 — solenoide de largo l, sección A y N vueltas con corriente I, con
% un núcleo de hierro dulce (mu) que ocupa un largo x; el resto es vacío.
% A corriente constante, F_x = dU/dx = (I^2 N^2 A / 2 l^2)(mu - mu0) > 0:
% el núcleo es atraído hacia adentro.
\begin{tikzpicture}[scale=1]

  \fill[figgray!8] (0,-0.8) rectangle (7,0.8);
  % nucleo
  \fill[figamber!35] (3.8,-0.62) rectangle (8.6,0.62);
  \draw[figamber, line width=0.9pt] (3.8,-0.62) rectangle (8.6,0.62);
  \node[figlbl] at (8.0,0) {$\mu$};
  \node[figlbl] at (1.9,0) {$\mu_0$};
  % solenoide
  \draw[figline] (0,0.8) -- (7,0.8) (0,-0.8) -- (7,-0.8);
  \foreach \xx in {0.2,0.6,...,6.8} {
    \draw[figblue, line width=1pt] (\xx,-0.86) -- (\xx+0.14,0.86);
  }
  \node[figlbl, text=figblue, anchor=east] at (-0.15,0.55) {$N,\ I$};

  % cotas
  \draw[figgray, line width=0.6pt, {Latex[length=1.6mm]}-{Latex[length=1.6mm]}] (0,-1.2) -- (7,-1.2);
  \node[figlbl, anchor=north] at (3.5,-1.25) {$\ell$};
  \draw[figgray, line width=0.6pt, {Latex[length=1.6mm]}-{Latex[length=1.6mm]}] (3.8,1.15) -- (7,1.15);
  \node[figlbl, anchor=south] at (5.4,1.2) {$x$};

  % fuerza: hacia x creciente, es decir, hacia adentro del solenoide
  \draw[figvecres] (10.0,0) -- (8.75,0);
  \node[figlbl, text=figred, anchor=south] at (9.4,0.1) {$F_x$};

\end{tikzpicture}
